\chapter{ΘΕΩΡΗΤΙΚΟ ΥΠΟΒΡΑΘΡΟ}
\section{ΕΙΣΑΓΩΓΗ}
Τα σύγχρονα ΣΗΕ υφίστανται αρκετές αλλαγές λόγω της μαζικής διείσδυσης  ΑΠΕ, καθώς και της σταδιακής απόσυρσης συμβατικών θερμικών μονάδων παραγωγής. Ως φυσικό επακόλουθο, μειώνεται η συνολική φυσική αδράνεια του συστήματος και μεταβάλλεται η δυναμική του συμπεριφορά, καθιστώντας το πιο ευάλωτο σε διαταραχές ισοζυγίου παραγωγής-ζήτησης. Στο πλαίσιο αυτό, η παρακολούθηση της δυναμικής ευστάθειας σε πραγματικό χρόνο αποτελεί επιτακτική ανάγκη. \cite{kundur1994,zeng2020inertia}
Η αξιόπιστη εποπτεία του συστήματος βασίζεται στη λήψη δεδομένων από PMUs, τοποθετημένες σε κομβικούς ζυγούς του δικτύου. Με τη χρήση τεχνικών ταυτοποίησης συστημάτων (System Identification), πραγματοποιείται εξαγωγή τοπικών εκτιμήσεων για τις ιδιοτιμές των δυναμικών αποκρίσεων. Οι συγκεκριμένες εκτιμήσεις, όμως, υπόκεινται σε θόρυβο μετρήσεων, τοπικές διακυμάνσεις και αριθμητικές αβεβαιότητες. Για την αντιμετώπιση των προαναφερθέντων περιορισμών, επιστρατεύονται αλγόριθμοι συσταδοποίησης (clustering), οι οποίοι φιλτράρουν και  αξιολογούν τα εκάστοτε δεδομένα, επιτρέποντας τον ακριβή προσδιορισμό των καθολικών ιδιοτιμών του συστήματος. \cite{estimation2021,schumacher2022clusteredr3ls}
Οι μεθοδολογίες, τα μαθηματικά μοντέλα και οι αλγόριθμοι που περιγράφονται στο θεωρητικό υπόβαθρο αποτελούν θεμελιώδη εργαλεία που έχουν προταθεί στη διεθνή βιβλιογραφία. Η εκτενής αναφορά και επεξήγησή τους κρίνεται απαραίτητη για την καλύτερη κατανόηση του αντικειμένου της συγκεκριμένης διπλωματικής.

\section{ΔΥΝΑΜΙΚΗ ΣΥΜΠΕΡΙΦΟΡΑ ΚΑΙ ΕΚΤΙΜΗΣΗ ΙΔΙΟΤΙΜΩΝ ΣΗΕ}
\subsection{ΗΛΕΚΤΡΟΜΗΧΑΝΙΚΕΣ ΤΑΛΑΝΤΩΣΕΙΣ ΚΑΙ ΔΥΝΑΜΙΚΗ ΕΥΣΤΑΘΕΙΑ}
Η ευστάθεια συχνότητας και η εμφάνιση μικρών διαταραχών συνδέονται άμεσα με την ικανότητα του ΣΗΕ να διατηρεί τη διασυνδεδεμένη λειτουργία του μετά από απότομες μεταβολές φορτίου και παραγωγής. Έπειτα από μια διαταραχή, προκαλούνται ηλεκτρομηχανικές ταλαντώσεις μεταξύ των περιστρεφόμενων μαζών των σύγχρονων γεννητριών. Οι ταλαντώσεις αυτές διακρίνονται σε:

\begin{itemize}
    \item \textbf{Intra-area Modes:} Αφορούν ταλαντώσεις μιας γεννήτριας ως προς το υπόλοιπο σύστημα, με συχνότητες συνήθως στο εύρος $0.8~\mathrm{Hz}$ έως $2.0~\mathrm{Hz}$.

    \item \textbf{Inter-area Modes:} Αφορούν ταλαντώσεις ομάδων γεννητριών μιας περιοχής ως προς γεννήτριες άλλης περιοχής μέσω ασθενών γραμμών μεταφοράς, με συχνότητες στο εύρος $0.1~\mathrm{Hz}$ έως $0.8~\mathrm{Hz}$.
\end{itemize}

Σε περίπτωση ανεπαρκούς απόσβεσης, οι ταλαντώσεις μπορεί να ενισχυθούν, οδηγώντας σε μη επιτρεπτούς ρυθμούς μεταβολής της συχνότητας (\textit{Rate of Change of Frequency} -- ROCOF). Παράλληλα, είναι πιθανό να προκληθούν αποκλίσεις τάσης και συχνότητας και, στη χειρότερη περίπτωση, διαχωρισμός του δικτύου και γενικευμένη αλληλουχία σφαλμάτων (\textit{blackout}). \cite{kundur1994,estimation2021}

\subsection{WAMS ΚΑΙ ΜΟΝΑΔΕΣ PMU}
Τα Συστήματα Παρακολούθησης Ευρείας Περιοχής (WAMS -- \textit{Wide-Area Monitoring Systems}) αποτελούν ένα από τα πιο τεχνολογικά εξελιγμένα συστήματα εποπτείας στα ΣΗΕ. Υλοποιήθηκαν για να βελτιώσουν τα υπάρχοντα συστήματα SCADA, προσφέροντας πιο αποτελεσματική παρακολούθηση της δυναμικής συμπεριφοράς του δικτύου σε ευρεία γεωγραφική έκταση και σε πραγματικό χρόνο. \cite{estimation2021,modal2018}

Η αρχιτεκτονική ενός συστήματος WAMS περιλαμβάνει:

\begin{itemize}
    \item \textbf{Μονάδες Μέτρησης Φασιθετών (PMUs):} Τοποθετούνται σε κρίσιμους υποσταθμούς για τη λήψη των πρωτογενών μετρήσεων.

    \item \textbf{Μονάδες Συγκέντρωσης Δεδομένων Φασιθετών (PDCs -- \textit{Phasor Data Concentrators}):} Συλλέγουν, ευθυγραμμίζουν χρονικά και αποθηκεύουν τα δεδομένα από πολλές PMUs.

    \item \textbf{Εφαρμογές Κέντρου Ελέγχου:} Επεξεργάζονται τα συγχρονισμένα δεδομένα για την εκτίμηση ευστάθειας, την ανίχνευση ταλαντώσεων και την υποστήριξη αποφάσεων.
\end{itemize}

Οι Μονάδες Μέτρησης Φασιθετών (\textit{Phasor Measurement Units} -- PMUs) είναι προηγμένες ηλεκτρονικές συσκευές που υπολογίζουν συγχρονισμένους φασιθέτες για μεγέθη όπως:

\begin{itemize}
    \item Το μέτρο ($V$, $I$) και τη γωνία φάσης τάσεων και ρευμάτων.

    \item Τη συχνότητα του δικτύου ($f$).

    \item Τον ρυθμό μεταβολής της συχνότητας (ROCOF).
\end{itemize}

Μέσω μιας προσέγγισης βασισμένης αποκλειστικά σε δεδομένα (\textit{Model-Free / Data-Driven approach}), η εκτίμηση των ιδιοτιμών του συστήματος, δηλαδή της συχνότητας ταλάντωσης και του συντελεστή απόσβεσης, πραγματοποιείται απευθείας από τις συγχρονισμένες χρονοσειρές. Παρακάμπτεται, έτσι, η ανάγκη για λεπτομερή γνώση των παραμέτρων ή της τοπολογίας του δικτύου, καθώς η πληροφορία εξάγεται απευθείας από την πραγματική απόκριση του συστήματος. \cite{modal2018,estimation2021}

Για την επεξεργασία των σημάτων σε πραγματικό χρόνο χρησιμοποιούνται προηγμένοι αλγόριθμοι, όπως οι μέθοδοι Prony, Matrix Pencil, Subspace Identification (SSI) και Dynamic Mode Decomposition (DMD). Οι προαναφερθείσες τεχνικές αναλύουν συνεχώς τις δυναμικές αποκρίσεις του δικτύου, επιτρέποντας τον έγκαιρο εντοπισμό ασθενώς αποσβεννύμενων ταλαντώσεων. Η διαρκής παρακολούθηση προσφέρει κρίσιμη πληροφορία στα Κέντρα Ελέγχου Ενέργειας, δίνοντας τη δυνατότητα προληπτικών παρεμβάσεων πριν οι ταλαντώσεις οδηγήσουν σε αστάθεια και σε γενικευμένη κατάρρευση του ηλεκτρικού δικτύου. \cite{modal2018,hua1990matrix,n4sid1994}

\subsection{ΜΑΘΗΜΑΤΙΚΗ ΑΝΑΠΑΡΑΣΤΑΣΗ ΣΥΣΤΗΜΑΤΩΝ}

Ένα γραμμικοποιημένο χρονικά αμετάβλητο σύστημα συνεχούς χρόνου γύρω από ένα σημείο ισορροπίας περιγράφεται στο χώρο από τις εξισώσεις:

\[
\frac{dx(t)}{dt} = A \cdot x(t) + B \cdot u(t)
\tag{2.1}
\label{eq:state-equation}
\]

\[
y(t) = C \cdot x(t) + D \cdot u(t)
\tag{2.2}
\label{eq:output-equation}
\]

όπου $x(t) \in \mathbb{R}^{n}$ είναι το διάνυσμα κατάστασης, $u(t) \in \mathbb{R}^{m}$ το διάνυσμα εισόδου και $y(t) \in \mathbb{R}^{p}$ το διάνυσμα εξόδου, το οποίο αντιστοιχεί στις διαθέσιμες μετρήσεις PMU. Οι πίνακες $A$, $B$, $C$ και $D$ είναι αντίστοιχα οι πίνακες κατάστασης, εισόδου, εξόδου και άμεσης διέλευσης του γραμμικοποιημένου μοντέλου.

Οι ιδιοτιμές $\lambda_i \in \mathbb{C}$ του πίνακα κατάστασης $A$ καθορίζουν τη δυναμική συμπεριφορά και την ευστάθεια του συστήματος. Για πραγματικούς πίνακες $A$, οι μη πραγματικές ιδιοτιμές εμφανίζονται σε συζυγή ζεύγη της μορφής:

\[
\lambda_i = \sigma_i \pm j \cdot \omega_i
\tag{2.3}
\label{eq:eigenvalue-pair}
\]

Για κάθε συζυγές ζεύγος, η συχνότητα ταλάντωσης $f_i$ και ο λόγος απόσβεσης $\zeta_i$ προκύπτουν από τα πραγματικά και φανταστικά μέρη της ιδιοτιμής ως εξής:

\[
f_i = \frac{\omega_i}{2\pi}
\tag{2.4}
\label{eq:oscillation-frequency}
\]

\[
\zeta_i = \frac{-\sigma_i}{\sqrt{\sigma_i^2 + \omega_i^2}}
\tag{2.5}
\label{eq:damping-ratio}
\]

Για ασυμπτωτική ευστάθεια του γραμμικοποιημένου συνεχούς συστήματος απαιτείται όλες οι ιδιοτιμές να έχουν αρνητικό πραγματικό μέρος, δηλαδή $\sigma_i < 0$. Όσο μικρότερος είναι ο λόγος απόσβεσης, τόσο ασθενέστερα αποσβεννύεται ο αντίστοιχος τρόπος ταλάντωσης· στην πρακτική παρακολούθηση των ΣΗΕ, τιμές $\zeta_i$ της τάξης του $3\%$--$5\%$ ή χαμηλότερες αξιολογούνται συνήθως ως ενδείξεις ασθενούς απόσβεσης και απαιτούν περαιτέρω διερεύνηση. \cite{kundur1994,pai1989}

\section{ΤΕΧΝΙΚΕΣ ΤΑΥΤΟΠΟΙΗΣΗΣ ΣΥΣΤΗΜΑΤΩΝ}

Οι τεχνικές ταυτοποίησης συστημάτων (system identification) αποτελούν τον θεμέλιο λίθο της προσέγγισης που βασίζεται σε δεδομένα (data-driven approach) για τη δυναμική ανάλυση των ΣΗΕ. Στόχος των τεχνικών είναι η εξαγωγή των δυναμικών χαρακτηριστικών του συστήματος (όπως οι ιδιοτιμές, οι συχνότητες ταλάντωσης, οι συντελεστές απόσβεσης και οι ρυθμοί ταλάντωσης) απευθείας από τις διαθέσιμες μετρήσεις, χωρίς τη χρήση αναλυτικών εξισώσεων. Ανάλογα με το πεδίο στο οποίο πραγματοποιείται η επεξεργασία των εκάστοτε δεδομένων, οι μέθοδοι ταυτοποίησης κατατάσσονται σε δύο κύριες κατηγορίες:

\begin{itemize}
    \item \textbf{Μέθοδοι Πεδίου Χρόνου (Time-Domain Methods)} εφαρμόζονται απευθείας στις χρονοσειρές των μετρούμενων αποκρίσεων του συστήματος, όπως καταγράφονται από τις μονάδες PMU είτε μετά από μια διαταραχή (ringdown data) είτε υπό συνθήκες φυσιολογικής λειτουργίας (ambient noise). Το κύριο πλεονέκτημα των μεθόδων πεδίου χρόνου έγκειται στην άμεση εφαρμογή τους σε πραγματικό χρόνο χωρίς να έχει προηγηθεί επεξεργασία των δεδομένων.

    \begin{itemize}
        \item Prony και Matrix Pencil (MP) βασίζονται στην προσέγγιση του σήματος ως άθροισμα μιγαδικών εκθετικών όρων, προσφέροντας υψηλή ακρίβεια στον υπολογισμό της συχνότητας και της απόσβεσης. \cite{hua1990matrix,modal2018}

        \item Eigensystem Realization Algorithm (ERA) και οι μέθοδοι αναγνώρισης υποχώρου (Subspace Identification, π.χ. N4SID) κατασκευάζουν ισοδύναμα μοντέλα χώρου καταστάσεων μέσω αλγεβρικών τεχνικών (όπως η Singular Value Decomposition - SVD), παρουσιάζοντας μεγάλη ανθεκτικότητα. \cite{juang1985ERA,n4sid1994}

        \item Prediction Error Method (PEM) στηρίζεται σε παραμετρικά στατιστικά μοντέλα (ARX, ARMAX) και στη βελτιστοποίηση των σφαλμάτων πρόβλεψης.
    \end{itemize}

    \item \textbf{Μέθοδοι Πεδίου Συχνότητας (Frequency-Domain Methods)} επεξεργάζονται το φάσμα συχνοτήτων των σημάτων (μέσω του μετασχηματισμού Fourier) ή τις μετρούμενες συναρτήσεις μεταφοράς και απόκρισης συχνότητας (Frequency Response Functions -- FRFs). Οι συγκεκριμένοι μέθοδοι χαρακτηρίζονται για την την ικανότητά τους να χειρίζονται συστήματα ευρέος φάσματος συχνοτήτων. Βρίσκοντας εφαρμογή σε μοντελοποίηση γραμμών μεταφοράς, καλωδίων και μετασχηματιστών.

    \begin{itemize}
        \item Vector Fitting είναι ένας επαναληπτικός αλγόριθμος που προσαρμόζει ρητές συναρτήσεις (rational functions) σε δεδομένα απόκρισης συχνότητας, αντικαθιστώντας σταδιακά τους πόλους του συστήματος έως ότου επιτευχθεί η επιθυμητή σύγκλιση. \cite{gustavsen1999vectorfitting}
    \end{itemize}
\end{itemize}

Η επιλογή μεταξύ των δύο κατηγοριών εξαρτάται από τη φύση των διαθέσιμων δεδομένων και την επιθυμητή εφαρμογή. Οι μέθοδοι πεδίου χρόνου προτιμώνται για τη συνεχή, online παρακολούθηση και την άμεση ανίχνευση ασταθειών από τις ροές δεδομένων των PMU.~\cite{sfetkos2024measurement,schumacher2022clusteredr3ls} Αντιθέτως οι μέθοδοι πεδίου συχνότητας προσφέρουν εξαιρετική ακρίβεια κατά την offline ανάλυση, τη σύνθεση ισοδύναμων δικτύων και τη μελέτη ευρέων περιοχών συχνοτήτων. \cite{modal2018, gustavsen1999vectorfitting}

\section{ΤΕΧΝΙΚΕΣ ΣΥΣΤΑΔΟΠΟΙΗΣΗΣ}

\subsection{ΑΝΑΓΚΗ ΣΥΣΤΑΔΟΠΟΙΗΣΗΣ}

Στην επιτήρηση των σύγχρονων ΣΗΕ, όπου βρίσκονται κατανεμημένες PMUs, η τοπική εφαρμογή αλγορίθμων ταυτοποίησης παράγει ένα μεγάλο πλήθος εκτιμήσεων για τις ιδιοτιμές του συστήματος. Παρόλο που όλες οι μονάδες PMU καταγράφουν τη δυναμική συμπεριφορά του ίδιου του φυσικού δικτύου, οι τοπικά υπολογιζόμενες ιδιοτιμές δεν συγκλίνουν σε ένα μοναδικό σημείο. Αντιθέτως, εμφανίζουν αισθητή διασπορά στο μιγαδικό επίπεδο, καθιστώντας αναγκαία τη συνδυαστική επεξεργασία των δεδομένων σε επίπεδο χώρου.

Η προαναφερθείσα διασπορά οφείλεται κυρίως σε τρεις παράγοντες:

\begin{itemize}
    \item \textbf{Θόρυβος:} Η παρουσία θορύβου στις μετρήσεις των αισθητήρων, καθώς και οι καθυστερήσεις και οι διακυμάνσεις στα κανάλια επικοινωνίας των PMUs, εισάγουν τυχαίες διαταραχές στις χρονοσειρές.

    \item Ο προσδιορισμός της χοροθέτησης κάθε τρόπου ταλάντωσης ποικίλλει σημαντικά ανάλογα με τη θέση του ζυγού στο δίκτυο. Για παράδειγμα, ένας ρυθμός ταλάντωσης μπορεί να είναι έντονα ορατός στις μετρήσεις ενός σταθμού παραγωγής, αλλά σχεδόν αφανής σε ένα απομακρυσμένο κέντρο φορτίου, επηρεάζοντας την ακρίβεια της πρώτης εμφάνισης της ταλάντωσης.

    \item \textbf{Μαθηματικοί / Ψευδείς Ρυθμοί:} Προκειμένου να αποφευχθεί η υποτίμηση της δυναμικής τάξης του συστήματος, οι αλγόριθμοι ταυτοποίησης (όπως οι Prony, SSI, ERA) χρησιμοποιούν υψηλή τάξη μοντέλου. Η υπερεκτίμηση αυτή οδηγεί στην εμφάνιση ψευδών ιδιοτιμών, οι οποίες αποτελούν αμιγώς αριθμητικά παράγωγα χωρίς φυσική υπόσταση στο δίκτυο.
\end{itemize}

Για την αντιμετώπιση των συγκεκριμένων προκλήσεων, κρίνεται επιβεβλημένη η εφαρμογή αλγορίθμων συσταδοποίησης, όπως οι K-Means, DBSCAN, OPTICS. Τέτοιοι αλγόριθμοι επεξεργάζονται το συνολικό πλήθος των τοπικών εκτιμήσεων στο μιγαδικό επίπεδο και εντοπίζουν τις περιοχές υψηλής πυκνότητας σημείων, οι οποίες αντιστοιχούν στους πραγματικούς φυσικούς τρόπους ταλάντωσης του συστήματος. Ταυτόχρονα, οι μεμονωμένες εκτιμήσεις που βρίσκονται σε απόσταση από τις κύριες συστάδες αναγνωρίζονται ως παράγωγα θορύβου ή ψευδείς ρυθμοί και απορρίπτονται αποτελεσματικά.

Η τελική εξαγωγή του κέντρου βάρους κάθε συστάδας παρέχει μια υψηλής ακρίβειας εκτίμηση για την πραγματική ιδιοτιμή του συστήματος. Με αποτέλεσμα να μπορεί να προσδιοριστεί με βεβαιότητα η πραγματική συχνότητα ταλάντωσης και ο κρίσιμος συντελεστής απόσβεσης. Ως φυσικό επακόλουθο, η τοπικά κατακερματισμένη και θορυβώδης πληροφορία μετατρέπεται σε αξιόπιστη, επιτρέποντας στα Κέντρα Ελέγχου Ενέργειας και στα αυτόματα συστήματα ελέγχου να παρακολουθούν την ευστάθεια του δικτύου σε πραγματικό χρόνο. \cite{estimation2021,schumacher2022clusteredr3ls}

\subsection{ΑΛΓΟΡΙΘΜΟΣ K-MEANS}

Ο αλγόριθμος k-means αποτελεί μία από τις πιο διαδεδομένες μεθόδους συσταδοποίησης. Στόχος του είναι ο διαμερισμός $N$ σημείων δεδομένων $x_i \in \mathbb{R}^d$ σε $k$ προκαθορισμένες συστάδες $C = \{C_1, C_2, ..., C_k\}$, ελαχιστοποιώντας το άθροισμα των τετραγωνικών αποστάσεων (Sum of Squared Errors -- SSE) από τα αντίστοιχα κέντρα συστάδων ($\mu_j$):

\[
SSE = \sum_{j=1}^{k} \sum_{x_i \in C_j} \left\| x_i - \mu_j \right\|^2
\]

Η διαδικασία εκτέλεσης του αλγορίθμου περιλαμβάνει τα εξής επαναληπτικά βήματα:

\begin{enumerate}
    \item Αρχικοποίηση: Επιλογή $k$ αρχικών τυχαίων κέντρων.
    
    \item Ανάθεση: Κάθε σημείο δεδομένων ανατίθεται στη συστάδα με το πλησιέστερο κέντρο βάσει Ευκλείδειας απόστασης.
    
    \item Ενημέρωση: Υπολογισμός των νέων κέντρων ως ο μέσος όρος των σημείων που ανήκουν σε κάθε συστάδα.
    
    \item Τερματισμός: Επανάληψη των βημάτων 2 και 3 έως ότου μεταβληθούν ελάχιστα τα κέντρα ή επιτευχθεί ο μέγιστος αριθμός επαναλήψεων.
\end{enumerate}

Ο αλγόριθμος συσταδοποίησης K-Means διακρίνεται για τη σημαντική υπολογιστική του ικανότητα και την υψηλή ταχύτητα εκτέλεσης, καθιστώντας τον μία από τις πιο δημοφιλείς επιλογές στην επεξεργασία δεδομένων. Η χαμηλή υπολογιστική του πολυπλοκότητα, επιτρέπει την ταχεία εκτέλεση και την αποτελεσματική διαχείριση μεγάλων συνόλων δεδομένων. \cite{macqueen1967some}

Παρά τα πλεονεκτήματα, ο αλγόριθμος παρουσιάζει ορισμένους σημαντικούς περιορισμούς που πρέπει να λαμβάνονται υπόψη. Κυριότερος εξ αυτών είναι η απαίτηση για εκ των προτέρων καθορισμό του αριθμού των συστάδων, μια τιμή που στην πράξη σπάνια είναι γνωστή εκ των προτέρων. Επιπλέον, ο K-Means παρουσιάζει μεγάλη ευαισθησία στην αρχική επιλογή των κέντρων, γεγονός που μπορεί να οδηγήσει  σε τοπικά ελάχιστα αντί της ολικής βέλτιστης λύσης. Οι συστάδες, επιπρόσθετα, έχουν σφαιρικό σχήμα και παρόμοιο μέγεθος, ενώ εμφανίζει ιδιαίτερη ευαισθησία στην παρουσία θορύβου και ακραίων τιμών (outliers).

Η επιλογή του αριθμού συστάδων $k$ μπορεί να υποστηριχθεί με τον δείκτη silhouette. Για κάθε σημείο υπολογίζεται η μέση απόστασή του από τα σημεία της ίδιας συστάδας και η μέση απόστασή του από την πλησιέστερη γειτονική συστάδα. Ο δείκτης silhouette λαμβάνει τιμές από $-1$ έως $1$. Τιμές κοντά στο $1$ υποδηλώνουν ότι το σημείο ανήκει σαφώς στη συστάδα του, τιμές κοντά στο $0$ δείχνουν επικάλυψη μεταξύ συστάδων, ενώ αρνητικές τιμές μπορεί να υποδηλώνουν εσφαλμένη ανάθεση. Για κάθε υποψήφια τιμή του $k$ υπολογίζεται ο μέσος δείκτης silhouette όλων των σημείων και επιλέγεται συνήθως η τιμή που τον μεγιστοποιεί. Με τον τρόπο αυτό, η επιλογή του αριθμού συστάδων γίνεται με πιο αντικειμενικό τρόπο. \cite{ROUSSEEUW198753}

\subsection{ΑΛΓΟΡΙΘΜΟΣ K-MEDOIDS}

Ο αλγόριθμος k-medoids είναι μέθοδος συσταδοποίησης παρόμοια με τον k-means, με τη βασική διαφορά ότι το αντιπροσωπευτικό στοιχείο κάθε συστάδας (medoid) είναι ένα πραγματικό σημείο του συνόλου δεδομένων και όχι ο μέσος όρος των σημείων της. Ο στόχος είναι η ελαχιστοποίηση του αθροίσματος των αποστάσεων κάθε σημείου από το medoid της συστάδας του.

Στην κλασική υλοποίηση PAM (Partitioning Around Medoids), επιλέγονται αρχικά $k$ medoids και στη συνέχεια εξετάζονται ανταλλαγές μεταξύ medoids και μη-medoid σημείων, έως ότου δεν μειώνεται άλλο το συνολικό κόστος. Η χρήση πραγματικών παρατηρήσεων ως κέντρων καθιστά τη μέθοδο πιο ανθεκτική σε ακραίες τιμές και ψευδείς ιδιοτιμές σε σχέση με τον k-means.

Ωστόσο, απαιτείται και εδώ εκ των προτέρων καθορισμός του αριθμού $k$, ενώ η υπολογιστική επιβάρυνση του PAM είναι συνήθως μεγαλύτερη. Για εκτιμήσεις ιδιοτιμών, ο k-medoids είναι χρήσιμος όταν επιδιώκεται ένα αντιπροσωπευτικό σημείο που να αντιστοιχεί σε πραγματική τοπική εκτίμηση και όχι σε τεχνητό μέσο όρο. \cite{kaufman2009finding}

\subsection{ΑΓΟΡΙΘΜΟΣ DBSCAN}

Ο αλγόριθμος DBSCAN είναι μια μέθοδος συσταδοποίησης που εντοπίζει συστάδες ως περιοχές υψηλής πυκνότητας σημείων που χωρίζονται από περιοχές χαμηλής πυκνότητας. Βασίζεται σε δύο θεμελιώδεις παραμέτρους:

\begin{itemize}
    \item $\varepsilon$ (eps): Η ακτίνα της γειτονιάς γύρω από ένα σημείο.

    \item MinPts: Ο ελάχιστος αριθμός σημείων που απαιτείται εντός της $\varepsilon$-γειτονιάς για να χαρακτηριστεί ένα σημείο ως κεντρικό.
\end{itemize}

Με βάση τις παραπάνω παραμέτρους, τα σημεία του χώρου ταξινομούνται σε τρεις κατηγορίες:

\begin{itemize}
    \item \textbf{Κεντρικά Σημεία (Core Points):} Σημεία που περιέχουν τουλάχιστον MinPts σημεία στην $\varepsilon$-γειτονιά τους (συμπεριλαμβανομένου του εαυτού τους).

    \item \textbf{Συνοριακά Σημεία (Border Points):} Σημεία που δεν είναι κεντρικά, αλλά βρίσκονται εντός της $\varepsilon$-γειτονιάς ενός κεντρικού σημείου.

    \item \textbf{Σημεία Θορύβου (Noise Points / Outliers):} Σημεία που δεν είναι ούτε κεντρικά ούτε συνοριακά.
\end{itemize}

Ο αλγόριθμος συνδέει άμεσα και έμμεσα προσπελάσιμα σημεία πυκνότητας για τη δημιουργία συστάδων αυθαίρετου σχήματος. Τα σημεία που δεν συνδέονται με καμία συστάδα ταξινομούνται αυτόματα ως θόρυβος.

Ο DBSCAN παρουσιάζει ιδιαίτερα πλεονεκτήματα, καθώς δεν απαιτεί τον εκ των προτέρων καθορισμό του αριθμού των συστάδων. Επιπλέον, διακρίνεται για την ικανότητά του να ανακαλύπτει συστάδες αυθαίρετου σχήματος, ενώ διαθέτει ενσωματωμένο μηχανισμό φιλτραρίσματος του θορύβου και των ψευδών ιδιοτιμών, απομονώνοντας αποτελεσματικά τα μεμονωμένα ανεπιθύμητα σημεία.

Από την άλλη πλευρά, ο αλγόριθμος εμφανίζει συγκεκριμένους περιορισμούς, με κυριότερο τη δυσκολία διαχωρισμού συστάδων που χαρακτηρίζονται από μεταβλητές πυκνότητες. Παράλληλα, η απόδοσή του εξαρτάται σε μεγάλο βαθμό από τη σωστή ρύθμιση των παραμέτρων του, δηλαδή της ακτίνας $\varepsilon$ και του ελάχιστου αριθμού απαιτούμενων σημείων MinPts. \cite{dbscan1996}

\subsection{ΑΓΟΡΙΘΜΟΣ OPTICS}

Ο αλγόριθμος OPTICS (Ordering Points To Identify the Clustering Structure) αποτελεί επέκταση του DBSCAN για δεδομένα στα οποία οι συστάδες ενδέχεται να έχουν διαφορετική πυκνότητα. Αντί για την παραγωγή μιας στατικής συσταδοποίησης για μια σταθερή τιμή $\varepsilon$, ο OPTICS δημιουργεί πρώτα μια ταξινομημένη ακολουθία των σημείων που αντιπροσωπεύει τη δομή πυκνότητάς τους.

Για κάθε σημείο $p$, εισάγονται δύο βασικές έννοιες:

\begin{itemize}
    \item \textbf{Απόσταση Πυρήνα (Core-distance):} Η τιμή αυτή ισούται με την απόσταση του $p$ από τον MinPts-οστό πλησιέστερο γείτονα του, υπό την προϋπόθεση ότι ο γείτονας βρίσκεται εντός της μέγιστης ακτίνας $\varepsilon$. Επομένως, μικρή core-distance υποδηλώνει ότι το σημείο βρίσκεται σε πυκνή περιοχή, ενώ για σημεία που δεν διαθέτουν επαρκή αριθμό γειτόνων η τιμή παραμένει αόριστη.

    \item \textbf{Απόσταση Προσπελασιμότητας (Reachability-distance):} εκφράζει το πόσο εύκολα ένα σημείο $p$ μπορεί να συνδεθεί με μία ήδη ανιχνευμένη πυκνή περιοχή μέσω ενός σημείου $o$. Ορίζεται ως η μέγιστη τιμή μεταξύ της core-distance του $o$ και της πραγματικής απόστασης μεταξύ των σημείων $o$ και $p$. Μια μικρή τιμή reachability-distance σημαίνει ότι το $p$ βρίσκεται κοντά σε πυκνή συστάδα, ενώ μια μεγάλη τιμή υποδεικνύει μετάβαση σε αραιότερη περιοχή ή σε διαφορετική συστάδα.
\end{itemize}

Με τη χρήση των προαναφερθέντων εννοιών, ο OPTICS παράγει το διάγραμμα προσπελασιμότητας (reachability plot), όπου οι κοιλάδες (valleys) αντιπροσωπεύουν συστάδες υψηλής πυκνότητας και οι κορυφές (peaks) διαχωρίζουν τις συστάδες ή υποδηλώνουν θόρυβο. Επιτρέποντας, έτσι, την εξαγωγή συστάδων διαφορετικής πυκνότητας χωρίς να απαιτείται επαναληπτική εκτέλεση του αλγορίθμου για πολλές τιμές της $\varepsilon$. \cite{optics1999}

\subsection{ΑΛΓΟΡΙΘΜΟΣ HDBSCAN}

Ο HDBSCAN (Hierarchical Density-Based Spatial Clustering of Applications with Noise) επεκτείνει τον DBSCAN, κατασκευάζοντας ιεραρχία συστάδων πυκνότητας αντί να βασίζεται σε μία μόνο τιμή της ακτίνας $\varepsilon$. Η μέθοδος χρησιμοποιεί την απόσταση αμοιβαίας προσπελασιμότητας (mutual reachability distance), η οποία ενσωματώνει την τοπική πυκνότητα των σημείων, και δημιουργεί ένα ελάχιστο δέντρο κάλυψης του συνόλου δεδομένων.

Η ιεραρχία συμπυκνώνεται με βάση την παράμετρο ελάχιστου μεγέθους συστάδας (min\_cluster\_size) και επιλέγονται οι πιο σταθερές συστάδες. Έτσι, ο αλγόριθμος μπορεί να εντοπίζει συστάδες διαφορετικής πυκνότητας και να χαρακτηρίζει σημεία χαμηλής σταθερότητας ως θόρυβο. Η παράμετρος min\_samples μπορεί να χρησιμοποιηθεί για αυστηρότερο ή ηπιότερο φιλτράρισμα θορύβου.

Στο πρόβλημα της ομαδοποίησης ιδιοτιμών, ο HDBSCAN είναι ιδιαίτερα κατάλληλος όταν οι πραγματικοί ρυθμοί ταλάντωσης εμφανίζουν άνισο πλήθος τοπικών εκτιμήσεων ή διαφορετική διασπορά. Δεν απαιτεί προκαθορισμένο αριθμό συστάδων, αλλά η επιλογή των παραμέτρων του πρέπει να συνδέεται με το πλήθος των διαθέσιμων PMUs και το αναμενόμενο επίπεδο επαναληψιμότητας των φυσικών ρυθμών. \cite{hdbscan2015}

\subsection{ΑΛΓΟΡΙΘΜΟΣ GAUSSIAN MIXTURE MODELS}

Τα Gaussian Mixture Models (GMMs) περιγράφουν την κατανομή των δεδομένων ως σταθμισμένο άθροισμα πολυδιάστατων κανονικών κατανομών. Κάθε συστάδα χαρακτηρίζεται από ένα μέσο διάνυσμα, έναν πίνακα συνδιακύμανσης και ένα βάρος μίξης. Σε αντίθεση με τον k-means, η ανάθεση ενός σημείου σε συστάδα είναι πιθανοτική (soft clustering) και όχι αποκλειστική.

Οι παράμετροι του μοντέλου εκτιμώνται συνήθως με τον αλγόριθμο Expectation--Maximization (EM), ο οποίος εναλλάσσει τον υπολογισμό των πιθανοτήτων συμμετοχής και την ενημέρωση των παραμέτρων των κανονικών κατανομών. Η δυνατότητα πλήρους πίνακα συνδιακύμανσης επιτρέπει την περιγραφή ελλειπτικών και προσανατολισμένων συστάδων στο επίπεδο των ιδιοτιμών.

Ο αριθμός συνιστωσών επιλέγεται, για παράδειγμα, με τα κριτήρια BIC ή AIC. Τα GMMs είναι χρήσιμα όταν ζητείται ποσοτική ένδειξη αβεβαιότητας για κάθε εκτίμηση ιδιοτιμής, αλλά επηρεάζονται από λανθασμένο αριθμό συνιστωσών και από ακραίες τιμές, οι οποίες συνήθως πρέπει να φιλτράρονται εκ των προτέρων. \cite{gmm2002,dempster1977em}

\subsection{ΑΛΓΟΡΙΘΜΟΣ AGGLOMERATIVE HIERARCHICAL CLUSTERING}

Η συσσωρευτική ιεραρχική συσταδοποίηση (Agglomerative Hierarchical Clustering) ακολουθεί προσέγγιση από τη βάση προς την κορυφή. Αρχικά, κάθε σημείο αποτελεί ξεχωριστή συστάδα και σε κάθε βήμα συγχωνεύονται οι δύο εγγύτερες συστάδες, μέχρι να επιτευχθεί ο επιθυμητός αριθμός συστάδων ή ένα προκαθορισμένο κατώφλι απόστασης.

Η απόσταση μεταξύ συστάδων καθορίζεται από κανόνα σύνδεσης (linkage), όπως single, complete, average ή Ward. Η μέθοδος Ward ελαχιστοποιεί την αύξηση της ενδοσυσταδιακής διασποράς και είναι συχνά κατάλληλη για συμπαγείς ομάδες εκτιμήσεων ιδιοτιμών. Το αποτέλεσμα απεικονίζεται σε δενδρόγραμμα, το οποίο επιτρέπει την επιλογή διαφορετικού επιπέδου διαμερισμού χωρίς επανεκτέλεση του αλγορίθμου.

Η μέθοδος δεν διαθέτει εγγενή μηχανισμό απόρριψης θορύβου και είναι ευαίσθητη στην επιλογή μετρικής και linkage. Επομένως, για εφαρμογές PMU συνιστάται να προηγείται βασικό φιλτράρισμα μη φυσικών εκτιμήσεων και να εφαρμόζεται κατάλληλη κανονικοποίηση των χαρακτηριστικών πριν από τον υπολογισμό των αποστάσεων. \cite{agglomerative1963}


\subsection{ΣΥΓΚΡΙΣΗ ΑΛΓΟΡΙΘΜΩΝ}

Ο Πίνακας \ref{tab:comparison-algorithms} συνοψίζει τα βασικά χαρακτηριστικά, τις απαιτήσεις και τους περιορισμούς των αλγορίθμων συσταδοποίησης που παρουσιάστηκαν.

\begin{landscape}
\thispagestyle{empty}
\AddToShipoutPictureFG*{%
  \AtPageLowerLeft{\put(570,423){\makebox(0,0){\rotatebox{90}{\thepage}}}}%
}
\begin{table}[H]
\centering
\caption{Συγκριτική παρουσίαση των αλγορίθμων συσταδοποίησης.}
\label{tab:comparison-algorithms}
\small
\setlength{\tabcolsep}{2pt}
\renewcommand{\arraystretch}{1.45}
\hspace*{-1.6cm}% Οπτική κεντράριση στη σελίδα landscape.
\resizebox{0.82\paperheight}{!}{%
\begin{tabular}{|p{3.1cm}|p{2.75cm}|p{2.75cm}|p{2.75cm}|p{2.75cm}|p{2.75cm}|p{2.75cm}|p{2.75cm}|}
\hline
\textbf{\scriptsize ΧΑΡΑΚΤΗΡΙΣΤΙΚΟ} & \textbf{\scriptsize K-MEANS} & \textbf{\scriptsize K-MEDOIDS} & \textbf{\scriptsize DBSCAN} & \textbf{\scriptsize OPTICS} & \textbf{\scriptsize HDBSCAN} & \textbf{\scriptsize GMM} & \textbf{\scriptsize AGGLOMERATIVE} \\ \hline
\textbf{Τύπος Αλγορίθμου} & Διαμερισμού & Διαμερισμού & Βασισμένος στην πυκνότητα & Πυκνότητας / ιεραρχικός & Ιεραρχικός πυκνότητας & Πιθανοτικός & Ιεραρχικός \\ \hline
\textbf{Απαιτούμενος αριθμός συστάδων} & Ναι & Ναι & Όχι & Όχι & Όχι & Ναι ή επιλογή με BIC/AIC & Όχι, αν οριστεί κατώφλι απόστασης \\ \hline
\textbf{Σχήμα συστάδων} & Σφαιρικό & Εξαρτάται από τη μετρική & Αυθαίρετο & Αυθαίρετο / μεταβλητή πυκνότητα & Αυθαίρετο / μεταβλητή πυκνότητα & Ελλειπτικό & Εξαρτάται από linkage \\ \hline
\textbf{Αντοχή σε θόρυβο / outliers} & Χαμηλή & Μέτρια έως υψηλή & Υψηλή & Υψηλή & Υψηλή & Χαμηλή χωρίς προφιλτράρισμα & Χαμηλή χωρίς προφιλτράρισμα \\ \hline
\textbf{Βασικές παράμετροι} & $k$ & $k$, μετρική απόστασης & eps, MinPts & MinPts, μέγιστη eps & min\_cluster\_size, min\_samples & Αριθμός συνιστωσών, τύπος συνδιακύμανσης & Linkage, μετρική, κατώφλι / αριθμός συστάδων \\ \hline
\textbf{Υπολογιστική απαίτηση} & Χαμηλή & Μεσαία έως υψηλή & Χαμηλή με χωρική ευρετηρίαση & Μεσαία & Μεσαία & Μεσαία & Μεσαία έως υψηλή \\ \hline
\end{tabular}
}
\end{table}
\end{landscape}

Η σύγκριση βασίζεται στις \cite{macqueen1967some,kaufman2009finding,dbscan1996,optics1999,hdbscan2015,gmm2002,agglomerative1963}.

\section{ΜΟΝΤΕΛΟΠΟΙΗΣΗ ΣΤΟΙΧΕΙΩΝ ΣΗΕ ΚΑΙ ΠΡΟΤΥΠΟ ΣΥΣΤΗΜΑ IEEE 39-BUS}

\subsection{ΜΟΝΑΔΕΣ ΠΑΡΑΓΩΓΗΣ}

Η δυναμική συμπεριφορά και η ευστάθεια των ΣΗΕ καθορίζονται άμεσα από τη σύνθεση και τα τεχνολογικά χαρακτηριστικά των εγκατεστημένων μονάδων παραγωγής. Στο συμβατικό ηλεκτρικό δίκτυο, πρωταγωνιστικό ρόλο της παραγωγής αποτελούν οι συμβατικές σύγχρονες γεννήτριες. Η δυναμική τους απόκριση περιγράφεται από τη θεμελιώδη εξίσωση ταλάντωσης, η οποία συνδέει τη μηχανική και την ηλεκτρική ισχύ μέσω της αποθηκευμένης κινητικής ενέργειας στους περιστρεφόμενους δρομείς τους. Επιπλέον, οι γεννήτριες πλαισιώνονται από κρίσιμα συστήματα ελέγχου, όπως οι αυτόματοι ρυθμιστές τάσης και οι ρυθμιστές στροφών. Η φυσική αδράνεια των περιστρεφόμενων μαζών προσφέρει αντίσταση στις απότομες μεταβολές της συχνότητας, απορροφώντας τις διαταραχές και σταθεροποιώντας τους ηλεκτρομηχανικούς ρυθμούς ταλάντωσης. \cite{kundur1994}

Η αυξημένη διείσδυση των ΑΠΕ άλλαξε ριζικά το τοπίο, με χαρακτηριστικό παράδειγμα τις ανεμογεννήτριες διπλά τροφοδοτούμενης επαγωγικής μηχανής (DFIG). Οι διατάξεις DFIG αξιοποιούν μια ασύγχρονη μηχανή με δακτυλίους, η οποία συνδέεται με το δίκτυο μέσω ενός μετατροπέα ισχύος στον δρομέα. Η παρουσία του μετατροπέα επιτρέπει την ηλεκτρική απόζευξη της ταχύτητας περιστροφής του δρομέα από τη συχνότητα του δικτύου. Ως αποτέλεσμα, η φυσική κινητική ενέργεια της ανεμογεννήτριας δεν συνεισφέρει άμεσα στη διατήρηση της συχνότητας. Παράλληλα, οι ελεγκτές του μετατροπέα εισάγουν νέους δυναμικούς τρόπους ταλάντωσης, οι οποίοι αλληλεπιδρούν με τους παραδοσιακούς ηλεκτρομηχανικούς ρυθμούς. \cite{zeng2020inertia}

Ακόμη πιο ουσιώδης είναι η επίδραση των μονάδων που συνδέονται στο δίκτυο αποκλειστικά μέσω μετατροπέων πηγής τάσης (Voltage Source Converters -- VSC), όπως τα φωτοβολταϊκά συστήματα, τα συστήματα αποθήκευσης ενέργειας (BESS) και οι σταθμοί μεταφοράς HVDC. Οι διατάξεις VSC στερούνται πλήρως περιστρεφόμενων μηχανικών μαζών και λειτουργούν μέσω γρήγορων βρόχων ελέγχου ρεύματος και τάσης. Η πλήρης απουσία φυσικής αδράνειας, σε συνδυασμό με την υψηλή ταχύτητα απόκρισης των ηλεκτρονικών ισχύος, μετατοπίζει τις ιδιοτιμές του συστήματος προς περιοχές υψηλότερων συχνοτήτων. Η συγκεκριμένη κατάσταση απαιτεί προηγμένη και λεπτομερή μοντελοποίηση, προκειμένου να αποτραπούν φαινόμενα ανεπιθύμητων αλληλεπιδράσεων ελέγχου και να διασφαλιστεί η ευστάθεια του συστήματος.

\subsection{ΜΟΝΤΕΛΑ ΦΟΡΤΙΩΝ}

Η μοντελοποίηση των φορτίων στα ΣΗΕ αποτελεί καθοριστικό παράγοντα για την ακριβή προσομοίωση της δυναμικής συμπεριφοράς και της ευστάθειας του δικτύου. Η εξάρτηση της απορροφούμενης ισχύος από την τάση και τη συχνότητα εκφράζεται ευρέως μέσω των στατικών μοντέλων ZIP, τα οποία συνδυάζουν τρεις βασικές συνιστώσες συμπεριφοράς φορτίου. \cite{kundur1994}

\begin{itemize}
    \item Η πρώτη συνιστώσα αφορά τα φορτία Σταθερής Σύνθετης Αντίστασης (Constant Impedance - Z), στα οποία η απορροφούμενη ισχύς μεταβάλλεται αναλογικά με το τετράγωνο της τάσης, σύμφωνα με τον λόγο
    $\left(\frac{V}{V_0}\right)^2$.

    \item Η δεύτερη συνιστώσα περιλαμβάνει τα φορτία Σταθερού Ρεύματος (Constant Current - I), όπου η ισχύς μεταβάλλεται γραμμικά με την τάση
    $\left(\frac{V}{V_0}\right)$.

    \item Η τρίτη συνιστώσα αντιπροσωπεύει τα φορτία Σταθερής Ισχύος (Constant Power - P), στα οποία η απορροφούμενη ισχύς διατηρείται σταθερή και παραμένει πλήρως ανεξάρτητη από τις διακυμάνσεις της τάσης.
\end{itemize}

Ο τύπος και η σύνθεση του φορτίου επηρεάζουν καταλυτικά τη δυναμική απόσβεση των ταλαντώσεων στο δίκτυο, διαμορφώνοντας τα όρια ευστάθειας μικρών διαταραχών. Ιδιαίτερα τα φορτία σταθερής ισχύος παρουσιάζουν μια τάση αποσταθεροποίησης των ηλεκτρομηχανικών ρυθμών ταλάντωσης. Αυτό συμβαίνει, διότι σε περίπτωση πτώσης της τάσης, το φορτίο αυξάνει αυτόματα την απορρόφηση ρεύματος προκειμένου να διατηρήσει την ισχύ του σταθερή. Επιβαρύνοντας, όμως, τις γραμμές μεταφοράς και εξασθενώντας την απόσβεση του συστήματος. \cite{kundur1994,sarkar2019impact}

Σε επίπεδο δυναμικής ανάλυσης, η συμπεριφορά εκδηλώνεται με τη μετατόπιση των κυρίαρχων ιδιοτιμών του συστήματος προς το δεξί μιγαδικό ημιεπίπεδο (αυξάνοντας το πραγματικό μέρος). Η συγκεκριμένη μετατόπιση μειώνει τον συντελεστή απόσβεσης και μπορεί να οδηγήσει σε ασταθείς ηλεκτρομηχανικές ταλαντώσεις αυξανόμενου πλάτους. Αναδεικνύοντας, έτσι, την ανάγκη για ακριβή εκτίμηση των παραμέτρων του μοντέλου ZIP στη σύγχρονη λειτουργία των ΣΗΕ. \cite{kundur1994,pai1989}

\subsection{ΠΡΟΤΥΠΟ ΣΥΣΤΗΜΑ IEEE 39-BUS}

Το πρότυπο σύστημα IEEE 39-bus (New England Power System) αποτελεί ευρέως αποδεκτό σύστημα αναφοράς για μελέτες δυναμικής ευστάθειας και ταλαντώσεων ευρείας περιοχής. Αποτελείται από 39 ζυγούς, 10 γεννήτριες (εκ των οποίων η μία αναπαριστά ισοδύναμο διασύνδεσης μεγάλης έκτασης), 19 φορτία και 46 γραμμές μεταφοράς/μετασχηματιστές. Το σύστημα εμφανίζει απόλυτα καθορισμένους τοπικούς και διασυνδετικούς ρυθμούς ταλάντωσης, καθιστώντας το ιδανικό για την αξιολόγηση και σύγκριση της ακρίβειας των αλγορίθμων συσταδοποίησης (k-means, DBSCAN, OPTICS) υπό διάφορα σενάρια διείσδυσης ΑΠΕ και τύπων φορτίου. \cite{digsilent39,estimation2021}
